\clearpage
\appendix
\section{Welfare bias in primitives}
\label{app:welfare-bias-primitives}
Recall $\varepsilon_r(p)\equiv p\,r'(p)/r(p)$ and
$\mu_p=1-r(p)-r'(p)(p+h)$.  Then
$r'(p)(p+h)=\varepsilon_r(p)\,r(p)\,(p+h)/p$, so
\begin{equation}
  \mu_p>0
  \;\iff\;
  \varepsilon_r(p)
  <
  \frac{p}{p+h}\,\frac{1-r(p)}{r(p)},
  \label{eq:app-mu-er}
\end{equation}

Result~\ref{res:welfare-bias} becomes:
\begin{enumerate}
\item[(i)] overestimation iff $\varepsilon_r(p)<0$;
\item[(ii)] underestimation iff $\varepsilon_r(p)>0$ and either $\mu_p>0$,
or both $\mu_p<0$ and $|\varepsilon_S|\le -\varepsilon_D$;
\item[(iii)] overestimation with wrong sign (predicted loss, actual gain)
iff $\mu_p<0$ and $|\varepsilon_S|>-\varepsilon_D$.
\end{enumerate}
\begin{proof}
By Result~\ref{res:welfare-bias}, overestimation with $0<\Lambda<1$ is
$r'<0$, underestimation is $r'>0$ with $\rho_c>0$, and overestimation with
wrong sign is $r'>0$ with $\rho_c<0$.
Since $\varepsilon_r$ has the same sign as $r'$, write these in
$\varepsilon_r$.  For~(i), $\varepsilon_r<0$ forces $\mu_p>1-r>0$ and hence
$\rho_c>0$, so $0<\Lambda<1$.
Result~\ref{res:negative-pt} translates the sign of $\rho_c$ into conditions
on $\mu_p$ and $|\varepsilon_S|\gtrless -\varepsilon_D$: with
$\varepsilon_r>0$ this is~(ii), and $\rho_c<0$ is~(iii).
\end{proof}
\begin{proof}[Proof of Corollary~\ref{cor:welfare-bias-buyer}]
The true change is $-Q\mu_p\Lambda$ and the no-return prediction is $-Q$, so
the no-return loss figure $Q$ understates the returns-adjusted figure
$Q\mu_p\Lambda$ if and only if $\mu_p\Lambda>1$ and overstates it if and only
if $\mu_p\Lambda<1$.  Write
$\mu_p=1-r(p)-r'(p)(p+h)$ and $\Lambda=1+r'(p)(p+h)\rho_c$.  Using
$\mu_p\rho_c=\varepsilon_S/(\varepsilon_S-\varepsilon_D)$,
\begin{align*}
  \mu_p\Lambda-1
  &=\mu_p\bigl(1+r'(p)(p+h)\rho_c\bigr)-1
  \\
  &=\mu_p-1+r'(p)(p+h)\,(\mu_p\rho_c)
  \\
  &=-r(p)-r'(p)(p+h)
    +r'(p)(p+h)\,\frac{\varepsilon_S}{\varepsilon_S-\varepsilon_D}
  \\
  &=-r(p)+r'(p)(p+h)\,\frac{\varepsilon_D}{\varepsilon_S-\varepsilon_D}
  \\
  &=-r(p)\Biggl(1-\varepsilon_r(p)\,\frac{p+h}{p}\,
       \frac{\varepsilon_D}{\varepsilon_S-\varepsilon_D}\Biggr).
\end{align*}
Hence $\mu_p\Lambda<1$ if and only if
$\varepsilon_r(p)\,\varepsilon_D/(\varepsilon_S-\varepsilon_D)<p/(p+h)$.
Multiplying by $-1$ gives
\[
  \varepsilon_r(p)\,
  \frac{-\varepsilon_D}{\varepsilon_S-\varepsilon_D}
  >
  -\frac{p}{p+h},
\]
which is \eqref{eq:V-buyer-cut}.  The reverse strict inequality is
$\mu_p\Lambda>1$.
\end{proof}

\section{Tax partial of the return rate}
\label{app:rtau}

The average return rate is $\tilde{r}(p,\tau)=N/D$ with
\[
  N=\int_p^\infty \phi(v,p-\tau)\,dF(v;\tau),
  \qquad
  D=1-F(p;\tau).
\]
Differentiating at fixed $p$,
\[
  \tilde{r}_{\tau}
  =\frac{N_{\tau}D-ND_{\tau}}{D^{2}}
  =\frac{N_{\tau}}{D}-\tilde{r}\,\frac{D_{\tau}}{D}.
\]
At $\tau=0$, $D=D(p)$ and $\tilde{r}=r(p)$.  Among remaining buyers, recovery
falls one-for-one, which contributes
$-\int_p^\infty \partial\phi(v,p)/\partial p\,dF(v)$ to $N_{\tau}$.  Buyer mass
changes at rate $D_{\tau}\le 0$.  The types who leave are those at $v=p$.  Over
an interval $d\tau$ they have measure $-D_{\tau}\,d\tau$ and cell return
probability $\phi(p,p)$, so $N$ changes by
\[
  -\phi(p,p)\,(-D_{\tau})\,d\tau
  =\phi(p,p)\,D_{\tau}\,d\tau.
\]
Collecting the two pieces,
\[
  N_{\tau}
  =-\int_p^\infty \frac{\partial\phi(v,p)}{\partial p}\,dF(v)
  +\phi(p,p)\,D_{\tau}.
\]
Substituting yields
\begin{equation}
  \tilde{r}_{\tau}
  =-\frac{1}{D(p)}\int_p^\infty \frac{\partial\phi(v,p)}{\partial p}\,dF(v)
  +\frac{-D_{\tau}}{D(p)}\bigl(r(p)-\phi(p,p)\bigr).
  \label{eq:rtau-split}
\end{equation}
